\subsection{TENSOR PRODUCT OF FAMILIES OF MULTIMODULES}

Let \((\mathbf{G}_{\lambda})_{\lambda \in L}\) be a family of \(\mathbf{Z}\) -modules; a mapping \(u\) of the set \(\mathbf{G} = \prod_{\lambda \in L} \mathbf{G}_{\lambda}\) into a \(\mathbf{Z}\) -module is said to be multiadditive (or \(\mathbf{Z}\) -multilinear) if \((x_{\lambda}) \mapsto u((x_{\lambda}))\) is additive with respect to each of the variables \(x_{\lambda}\) ; to be precise, this means that, for all \(\mu \in L\) and every element \((a_{\lambda}) \in \prod_{\lambda \neq \mu} \mathbf{G}_{\lambda}\) , canonically identifying \(\mathbf{G}\) with \(\mathbf{G}_{\mu} \times \prod_{\lambda \neq \mu} \mathbf{G}_{\lambda}\) ,

\[
u (x _ {\mu} + y _ {\mu}, (a _ {\lambda})) = u (x _ {\mu}, (a _ {\lambda})) + u (y _ {\mu}, (a _ {\lambda})) \quad \text { for } x _ {\mu}, y _ {\mu} \text { in } G _ {\mu}.\tag{24}
\]

This implies in particular that \(u((x_{\lambda})) = 0\) if one of the \(x_{\lambda}\) is zero.

We consider also the universal mapping problem where \(\Sigma\) is the species of Z-module structure and the \(\alpha\) -mappings are the multiadditive mappings of G into a Z-module. A solution is still obtained by considering the Z-module \(C = Z^{(G)}\) of formal linear combinations of elements of G with coefficients in Z and the sub-Z-module D of C generated by the elements of the form

\[
\left(x _ {\mu} + y _ {\mu}, \left(z _ {\lambda}\right) _ {\lambda \neq \mu}\right) - \left(x _ {\mu}, \left(z _ {\lambda}\right) _ {\lambda \neq \mu}\right) - \left(y _ {\mu}, \left(z _ {\lambda}\right) _ {\lambda \neq \mu}\right)
\]

where \(\mu \in \mathbf{L}\) , \(x_{\mu} \in \mathrm{G}_{\mu}\) , \(y_{\mu} \in \mathrm{G}_{\mu}\) and the \(z_{\lambda} \in \mathrm{G}_{\lambda}\) ( \(\lambda \neq \mu\) ) are arbitrary. The quotient \(\mathbf{Z}\) -module C/D is called the tensor product (over \(\mathbf{Z}\) ) of the family \((\mathrm{G}_{\lambda})_{\lambda \in L}\) of \(\mathbf{Z}\) -modules and is denoted by \(\bigotimes_{\lambda \in L} \mathrm{G}_{\lambda}\) ; for every element \((x_{\lambda})_{\lambda \in L}\) of G which is an element of the canonical basis of C, \(\bigotimes_{\lambda \in L} x_{\lambda}\) denotes the canonical image of this element in C/D. It follows from the above definitions that the mapping \(\phi: (x_{\lambda}) \mapsto \bigotimes_{\lambda \in L} x_{\lambda}\) of G into \(\bigotimes_{\lambda \in L} G_{\lambda}\) is Z-multilinear and that, for every Z-multilinear mapping \(f\) of G into a Z-module H, there exists one and only one Z-linear mapping \(g: \bigotimes_{\lambda \in L} G_{\lambda} \to H\) such that \(f = g \circ \phi\) ; the ordered pair \(\left( \bigotimes_{\lambda \in L} G_{\lambda}, \phi \right)\) thus resolves the universal mapping problem in question.

Let \((\mathbf{G}_{\lambda}^{\prime})_{\lambda \in L}\) be another family of \(\mathbf{Z}\) -modules and, for all \(\lambda \in L\) , let \(v_{\lambda}: \mathbf{G}_{\lambda} \to \mathbf{G}_{\lambda}^{\prime}\) be a \(\mathbf{Z}\) -linear mapping (in other words a homomorphism of commutative groups). Then the mapping

\[
(x _ {\lambda}) \mapsto \bigotimes_ {\lambda \in L} v _ {\lambda} (x _ {\lambda})
\]

of \(\mathbf{G}\) into \(\bigotimes_{\lambda \in L} \mathbf{G}_{\lambda}'\) is \(\mathbf{Z}\) -multilinear and hence defines canonically a \(\mathbf{Z}\) -linear mapping of \(\bigotimes_{\lambda \in L} \mathbf{G}_{\lambda}\) into \(\bigotimes_{\lambda \in L} \mathbf{G}_{\lambda}'\) denoted by \(\bigotimes_{\lambda \in L} v_{\lambda}\) and such that

\[
\left(\bigotimes_ {\lambda \in L} v _ {\lambda}\right) \left(\bigotimes_ {\lambda \in L} x _ {\lambda}\right) = \bigotimes_ {\lambda \in L} v _ {\lambda} (x _ {\lambda}).\tag{25}
\]

In particular, we consider, for some \(\mu \in \mathbf{L}\) , an endomorphism \(\theta\) of \(\mathbf{G}_{\mu}\) ; we denote by \(\tilde{\theta}\) the endomorphism of \(\bigotimes_{\lambda \in \mathbf{L}} \mathbf{G}_{\lambda}\) equal to \(\bigotimes_{\lambda \in \mathbf{L}} v_{\lambda}\) where \(v_{\mu} = \theta\) and \(v_{\lambda} = 1_{\mathbf{G}_{\lambda}}\) for \(\lambda \neq \mu\) .

Then we suppose given a set \(\Omega\) , a mapping

\[
c: \omega \mapsto (\rho (\omega), \sigma (\omega))
\]

of \(\Omega\) into \(\mathbf{L} \times \mathbf{L}\) and, for all \(\omega \in \Omega\) , an endomorphism \(p_{\omega}\) of \(\mathrm{G}_{\rho(\omega)}\) and an endomorphism \(q_{\omega}\) of \(\mathrm{G}_{\sigma(\omega)}\) ; there correspond to them two endomorphisms \(\tilde{p}_{\omega}\) and \(\tilde{q}_{\omega}\) of \(\mathrm{P} = \bigotimes_{\lambda \in \mathbf{L}} \mathrm{G}_{\lambda}\) . Let \(\mathbf{R}\) be the sub- \(\mathbf{Z}\) -module of \(\mathbf{P}\) generated by the union of the images of the endomorphisms \(\tilde{p}_{\omega} - \tilde{q}_{\omega}\) when \(\omega\) runs through \(\Omega\) . The quotient \(\mathbf{Z}\) -module \(\mathrm{P/R}\) is called the tensor product of the family \((\mathrm{G}_{\lambda})_{\lambda \in \mathbf{L}}\) relative to \(c, p, q\) and is denoted by \(\bigotimes_{(c, p, q)} \mathrm{G}_{\lambda}\) ; composing the canonical homomorphism \(\mathrm{P} \to \mathrm{P/R}\) with the mapping \(\phi: \mathrm{G} \to \bigotimes_{\lambda \in \mathbf{L}} \mathrm{G}_{\lambda}\) defined above, we obtain a \(\mathbf{Z}\) -multilinear mapping \(\phi_{(c, p, q)}: \mathrm{G} \to \bigotimes_{(c, p, q)} \mathrm{G}_{\lambda}\) and write \(\phi_{(c, p, q)}((x_{\lambda})) = \bigotimes_{(c, p, q)} x_{\lambda}\) or simply \(\bigotimes_{(c)} x_{\lambda}\) . The ordered pair consisting of \(\bigotimes_{(c, p, q)} x_{\lambda}\) and \(\phi_{(c, p, q)}\) resolves the following universal mapping problem: let \(\tilde{p}_{\omega}\) (resp. \(\tilde{q}_{\omega}\) ) denote the mapping

\[
\left(x _ {\rho (\omega)}, \left(x _ {\lambda}\right) _ {\lambda \neq \rho (\omega)}\right) \mapsto \left(p _ {\omega} \left(x _ {\rho (\omega)}\right), \left(x _ {\lambda}\right) _ {\lambda \neq \rho (\omega)}\right)
\]

\[
\left(\text { resp. } \left(x _ {\sigma (\omega)}, \left(x _ {\lambda}\right) _ {\lambda \neq \sigma (\omega)}\right) \mapsto \left(q _ {\omega} (x _ {\sigma (\omega)}), \left(x _ {\lambda}\right) _ {\lambda \neq \sigma (\omega)}\right)\right)
\]

of G into itself. Then \(\Sigma\) is taken to be the species of Z-module structure and the \(\alpha\) -mappings the \(\mathbf{Z}\) -multilinear mappings \(u\) of \(\mathbf{G}\) into a \(\mathbf{Z}\) -module which further satisfy the conditions

\[
u \circ \tilde {p} _ {\omega} = u \circ \tilde {q} _ {\omega}\tag{26}
\]

for all \(\omega \in \Omega\) . The proof is obvious from the above definitions.

This construction recovers in particular that of \(\mathbf{E} \otimes_{\mathbf{A}} \mathbf{F}\) described in no. 1: in this case we take \(\mathbf{L} = \{1, 2\}\) , \(\mathbf{G}_1 = \mathbf{E}\) , \(\mathbf{G}_2 = \mathbf{F}\) , \(\Omega = \mathbf{A}\) ; further, for all \(\omega \in \mathbf{A}\) , we must have \(\rho(\omega) = 1\) , \(\sigma(\omega) = 2\) , \(p_\omega\) is the endomorphism \(x \mapsto x\omega\) of the \(\mathbf{Z}\) -module \(\mathbf{E}\) and \(q_\omega\) the endomorphism \(y \mapsto \omega y\) of the \(\mathbf{Z}\) -module \(\mathbf{F}\) .

Let \((\mathbf{G}_{\lambda}^{\prime})_{\lambda \in L}\) be a second family of Z-modules; keeping the mapping c the same, suppose given for all \(\omega \in \Omega\) , an endomorphism \(p_{\omega}^{\prime}\) of \(\mathbf{G}_{\rho(\omega)}^{\prime}\) and an endomorphism \(q_{\omega}^{\prime}\) of \(\mathbf{G}_{\sigma(\omega)}^{\prime}\) . For all \(\lambda \in L\) , then let \(v_{\lambda}: G_{\lambda} \to G_{\lambda}^{\prime}\) be a Z-linear mapping such that, for all \(\omega \in \Omega\) ,

\[
v _ {\rho (\omega)} \circ p _ {\omega} = p _ {\omega} ^ {\prime} \circ v _ {\rho (\omega)} \quad \text { and } \quad v _ {\sigma (\omega)} \circ q _ {\omega} = q _ {\omega} ^ {\prime} \circ v _ {\sigma (\omega)}\tag{27}
\]

(in other words, for all \(\lambda \in \mathbf{L}\) , \(v_{\lambda}\) is a morphism for the laws of action on \(\mathbf{G}_{\lambda}\) (resp. \(\mathrm{G}_{\lambda}^{\prime}\) ) defined on the \(p_{\xi}\) and \(q_{\eta}\) (resp. \(p_{\xi}^{\prime}\) and \(q_{\eta}^{\prime}\) ) with \(\xi\) and \(\eta\) such that \(\rho(\xi) = \lambda\) and \(\sigma(\eta) = \lambda\) ). Then the mapping

\[
u: (x _ {\lambda}) \mapsto \bigotimes_ {(c, p', q')} v _ {\lambda} (x _ {\lambda})
\]

of G into \(\bigotimes_{(c, p', q')} \mathrm{G}_{\lambda}'\) is Z-multilinear and satisfies conditions (26) and hence defines a Z-linear mapping of \(\bigotimes_{(c, p, q)} \mathrm{G}_{\lambda}\) into \(\bigotimes_{(c, p', q')} \mathrm{G}_{\lambda}'\) , which we shall denote simply by \(\bigotimes_{(c)} v_{\lambda}\) if no confusion can arise.

We shall now give an ``associativity'' property for the general tensor products thus defined. Let \((\mathbf{L}_{i})_{1\leq i\leq n}\) be a finite partition of L; for every index i, let \(\Omega_{i}\) denote the subset of \(\Omega\) consisting of the elements such that \(\rho(\omega)\in\mathbf{L}_{i}\) and \(\sigma(\omega)\in\mathbf{L}_{i}\) ; clearly the \(\Omega_{i}\) are pairwise disjoint; we set \(\Omega^{\prime}=\Omega-\left(\bigcup_{i}\Omega_{i}\right)\) . For each index i, \(c^{(i)}\) will denote the mapping \(\omega\mapsto(\rho(\omega),\sigma(\omega))\) of \(\Omega_{i}\) into \(L_{i}\times L_{i}\) ; for \(\omega\in\Omega_{i}\) , we write \(p_{\omega}^{(i)}\) and \(q_{\omega}^{(i)}\) instead of \(p_{\omega}\) and \(q_{\omega}\) . Then for each i there is a ``partial'' tensor product

\[
\mathrm{F} _ {i} = \bigotimes_ {(c ^ {(i)}, p ^ {(i)}, q ^ {(i)})} \mathrm{G} _ {\lambda}.
\]

We shall further make the following ``permutability'' hypothesis:

\begin{enumerate}
\def\labelenumi{(\Alph{enumi})}
\setcounter{enumi}{15}
\tightlist
\item
  If \(\omega \in \Omega'\) , \(p_{\omega}\) (resp. \(q_{\omega}\) ) permutes with each of the endomorphisms \(p_{\xi}\) and \(q_{\eta}\) of \(\mathrm{G}_{\rho(\omega)}\) (resp. \(\mathrm{G}_{\sigma(\omega)}\) ) such that \(\xi \notin \Omega'\) , \(\eta \notin \Omega'\) and \(\rho(\omega) = \rho(\xi) = \sigma(\eta)\) (resp. \(\sigma(\omega) = \rho(\xi) = \sigma(\eta)\) ).
\end{enumerate}

For each \(\omega \in \Omega'\) , let \(i\) be the index such that \(\rho(\omega) \in \mathbf{L}_i\) ; then consider the family \((v_{\lambda})_{\lambda \in \mathbf{L}_i}\) where \(v_{\rho(\omega)} = p_{\omega}\) and \(v_{\lambda} = 1_{\mathbf{G}_{\lambda}}\) for \(\lambda \neq \rho(\omega)\) ; hypothesis (P) implies that the family \((v_{\lambda})\) satisfies conditions (27) (where \(p'\) and \(p\) must be replaced by \(p^{(i)}\) , \(q'\) and \(q\) by \(q^{(i)}\) , \(\omega\) by an element \(\xi\) running through \(\Omega_i\) ); thus an endomorphism \(\bigotimes_{(c^{(i)})} v_{\lambda} = r_{\omega}\) is derived of the \(\mathbf{Z}\) -module \(\mathbf{F}_i\) . Similarly, an endomorphism \(s_{\omega}\) is defined of the \(\mathbf{Z}\) -module \(\mathbf{F}_j\) starting with \(q_{\omega}, j\) being the index such that \(\sigma(\omega) \in \mathbf{L}_j\) ; finally let \(d(\omega) = (i,j)\) . Then we can define the tensor product \(\bigotimes_{(d,r,s)} \mathbf{F}_i\) and the corresponding canonical mapping

\[
\phi_ {(d, r, s)}: \prod_ {i = 1} ^ {n} \mathrm{F} _ {i} \rightarrow \bigotimes_ {(d, r, s)} \mathrm{F} _ {i}.
\]

On the other hand, for each i, the canonical mapping

\[
\psi_ {i} = \phi_ {(c ^ {(i)}, p ^ {(i)}, q ^ {(i)})}: \prod_ {\lambda \in L _ {i}} G _ {\lambda} \rightarrow F _ {i};
\]

using the associativity of the product of sets, a Z-multilinear mapping \(\psi = \phi_{(d,r,s)}\circ (\psi_i)\) of G into \(\bigotimes_{(d,r,s)}\mathrm{F}_i\) is derived. We show that the ordered pair \(\left(\bigotimes_{(d,r,s)}\mathrm{F}_i,\psi\right)\) is a solution of the same universal problem as \(\left(\bigotimes_{(c,p,q)}\mathrm{G}_{\lambda},\phi_{(c,p,q)}\right)\) , whence will follow the existence of a unique Z-module isomorphism

\[
\theta \colon \bigotimes_ {(c, p, q)} \mathrm{G} _ {\lambda} \rightarrow \bigotimes_ {(d, r, s)} \mathrm{F} _ {i}
\]

such that \(\psi = \theta \circ \phi_{(c,p,q)}\) (Set Theory, IV, § 3, no. 1).

By induction on \(n\) , the proof is reduced to the case \(n = 2\) ; for simplicity we write \(\mathbf{F}_1 \otimes_{(d)} \mathbf{F}_2\) and \(y_1 \otimes_{(d)} y_2\) instead of \(\bigotimes_{(d, r, s)} \mathbf{F}_i\) and \(\bigotimes_{(d, r, s)} y_i\) . Consider the mapping from \(G\) to \(\mathbf{F}_1 \otimes_{(d)} \mathbf{F}_2\)

\[
h \colon (x _ {\lambda}) \to \left(\bigotimes_ {(c ^ {(1)})} x _ {\lambda}\right) \underset {(d)} {\otimes} \left(\bigotimes_ {(c ^ {(2)})} x _ {\lambda}\right).
\]

It is obviously Z-multilinear; we show that it satisfies conditions (26) for all \(\omega\in\Omega\) . This is obvious if \(\omega\in\Omega_{1}\) or \(\omega\in\Omega_{2}\) ; otherwise, supposing, in order to fix the ideas, that \(\rho(\omega)\in\mathbf{L}_{1}\) and \(\sigma(\omega)\in\mathbf{L}_{2}\) , the values of \(h\circ\tilde{p}_{\omega}\) and \(h\circ\tilde{q}_{\omega}\) at \((x_{\lambda})\) are respectively

\[
\left(r _ {\omega} \left(\bigotimes_ {(c ^ {(1)})} x _ {\lambda}\right)\right) \underset {(d)} {\otimes} \left(\bigotimes_ {(c ^ {(2)})} x _ {\lambda}\right) \quad \text { and } \quad \left(\bigotimes_ {(c ^ {(1)})} x _ {\lambda}\right) \underset {(d)} {\otimes} \left(s _ {\omega} \left(\bigotimes_ {(c ^ {(2)})} x _ {\lambda}\right)\right)
\]

which are also equal by definition of \(\mathbf{F}_1\otimes_{(d)}\mathbf{F}_2\)

This being so, let \(u\) be a \(\mathbf{Z}\) -multilinear mapping of \(G\) into a \(\mathbf{Z}\) -module \(H\) , satisfying conditions (26); we shall define a \(\mathbf{Z}\) -linear mapping \(v: F_1 \otimes F_2 \to H\) such that \(u = v \circ h\) and that will prove our assertion (repeating the argument of no. 1, Corollary 1 to Proposition 1). For all \(z_2 = (x_\lambda)_{\lambda \in L_2}\) we consider the ``partial'' linear mapping of \(\prod_{\lambda \in L_1} G_\lambda\) into \(H\)

\[
u \left(\cdot , z _ {2}\right): \left(x _ {\lambda}\right) _ {\lambda \in L _ {1}} \mapsto u \left(\left(x _ {\lambda}\right) _ {\lambda \in L _ {1}}, z _ {2}\right) = u \left(\left(x _ {\lambda}\right) _ {\lambda \in L}\right).\tag{28}
\]

Clearly it is Z-multilinear and satisfies conditions (26) for \(\omega\in\Omega_{1}\) ; by definition there thus exists a Z-linear mapping \(y_{1}\mapsto w_{1}(y_{1},x_{2})\) of \(F_{1}\) into H such that

\[
w _ {1} \left(\bigotimes_ {(c ^ {(1)})} x _ {\lambda}, z _ {2}\right) = u \left(\left(x _ {\lambda}\right) _ {\lambda \in L _ {1}}, z _ {2}\right)\tag{29}
\]

We next consider the mapping

\[
u _ {2}: (x _ {\lambda}) _ {\lambda \in L _ {2}} \mapsto w _ {1} (\cdot , (x _ {\lambda}) _ {\lambda \in L _ {2}})
\]

of \(\prod_{\lambda \in L_2} G_\lambda\) into \(\mathrm{Hom}_{\mathbf{Z}}(\mathrm{F}_1, \mathrm{H})\) ; it is obviously \(\mathbf{Z}\) -multilinear and satisfies conditions (26) for \(\omega \in \Omega_2\) , by virtue of the hypothesis on \(u\) and relations (28) and (29) and taking account of the fact that the elements of the form \(\bigotimes_{(c^{(1)})} x_\lambda\) generate the \(\mathbf{Z}\) -module \(\mathrm{F}_1\) . Hence there is a \(\mathbf{Z}\) -linear mapping

\[
w _ {2}: \mathrm{F} _ {2} \rightarrow \operatorname{Hom} _ {\mathbf {Z}} (\mathrm{F} _ {1}, \mathrm{H})
\]

such that

\[
w _ {2} \left(\bigotimes_ {(c ^ {(2)})} x _ {\lambda}\right) = u _ {2} \left(\left(x _ {\lambda}\right) _ {\lambda \in L _ {2}}\right)
\]

or also

\[
\left(w _ {2} \left(\bigotimes_ {(c ^ {(2)})} x _ {\lambda}\right)\right) \left(\bigotimes_ {(c ^ {(1)})} x _ {\lambda}\right) = u \left(\left(x _ {\lambda}\right) _ {\lambda \in L}\right).\tag{30}
\]

We now consider, for \(y_{1} \in F_{1}, y_{2} \in F_{2}\) , the element of H

\[
w (y _ {1}, y _ {2}) = (w _ {2} (y _ {2})) (y _ {1}).\tag{31}
\]

Clearly \(w\) is a \(\mathbf{Z}\) -bilinear mapping of \(\mathbf{F}_1 \times \mathbf{F}_2\) into \(H\) . We show further that, for all \(\omega \in \Omega'\) , (assuming, to fix the ideas, that \(\rho(\omega) \in L_1\) and \(\sigma(\omega) \in L_2\) )

\[
w (r _ {\omega} (y _ {1}), y _ {2}) = w (y _ {1}, s _ {\omega} (y _ {2})).\tag{32}
\]

It suffices to verify this relation when \(y_{1}\) (resp. \(y_{2}\) ) is of the form \(\bigotimes_{(c^{(1)})} x_{\lambda}\) (resp. \(\bigotimes_{(c^{(2)})} x_{\lambda}\) ), since these elements generate the Z-module \(F_{1}\) (resp. \(F_{2}\) ). But by definition, \(r_{\omega}\left(\bigotimes_{(c^{(1)})} x_{\lambda}\right) = \bigotimes_{(c^{(1)})} x_{\lambda}'\) , where \(x_{\rho(\omega)}' = p_{\omega}(x_{\rho(\omega)})\) and \(x_{\lambda}' = x_{\lambda}\) for \(\lambda \neq \rho (\omega)\) in \(\mathbf{L}_1\) ; similarly \(s_{\omega}\left(\bigotimes_{(c^{(2)})}x_{\lambda}\right) = \bigotimes_{(c^{(2)})}x_{\lambda}^{\prime \prime}\) , where \(x_{\sigma (\omega)}^{\prime \prime} = q_{\omega}(x_{\sigma (\omega)})\) and \(x_{\lambda}^{\prime \prime} = x_{\lambda}\) for \(\lambda \neq \sigma (\omega)\) in \(\mathbf{L}_2\) ; using (30) and (31), relation (32) then follows from (26). Hence there exists a Z-linear mapping \(v\) of \(\mathrm{F}_1\otimes \mathrm{F}_2\) into H such that \(v(y_{1}\otimes y_{2}) = w(y_{1},y_{2})\) and it then follows from (30) and (31) that \(v\circ h = u\) .

The most important special case of the general tensor product defined above is the following: we start with a family \((\mathbf{A}_i)_{1\leqslant i\leqslant n - 1}\) of rings and a family \((\mathbf{E}_i)_{1\leqslant i\leqslant n}\) , where \(\mathbf{E}_1\) is a right \(\mathbf{A}_1\) -module, \(\mathbf{E}_n\) is a left \(\mathbf{A}_{n - 1}\) module and for \(2\leqslant i\leqslant n - 1\) , \(\mathbf{E}_i\) is an \((\mathbf{A}_{i - 1},\mathbf{A}_i)\) -bimodule. Then the above definition is applied as follows: L is the set \([1,n]\) , \(\mathrm{G}_i = \mathrm{E}_i\) , \(\Omega\) is the set the sum of the \(\mathbf{A}_i\) ( \(1\leqslant i\leqslant n - 1\) ). For \(\omega \in \mathbf{A}_i\) ( \(1\leqslant i\leqslant n - 1\) ), take \(\rho (\omega) = i\) , \(\sigma (\omega) = i + 1\) , \(p_{\omega}\) is the endomorphism \(x\mapsto x\omega\) of the Z-module \(\mathbf{E}_i\) and \(q_{\omega}\) the endomorphism \(y\mapsto \omega y\) of the Z-module \(\mathbf{E}_{i + 1}\) ; the corresponding tensor product is denoted by

\[
\mathbf {E} _ {1} \otimes_ {\mathbf {A} _ {1}} \mathbf {E} _ {2} \otimes_ {\mathbf {A} _ {2}} \mathbf {E} _ {3} \otimes \dots \otimes_ {\mathbf {A} _ {n - 2}} \mathbf {E} _ {n - 1} \otimes_ {\mathbf {A} _ {n - 1}} \mathbf {E} _ {n}\tag{33}
\]

(a notation where the \(A_{i}\) may occasionally be suppressed) and the elements \(\bigotimes_{(c,p,q)} x_{i}\) of this tensor product, for a family \((x_{i})\) such that \(x_{i} \in E_{i}\) for \(1 \leqslant i \leqslant n\) , may be written \(x_{1} \otimes x_{2} \otimes \cdots \otimes x_{n}\) if no confusion can arise; an analogous notation is used for a Z-linear mapping \(\bigotimes_{(c)} v_{i}\) . Hypothesis (P) holds for every partition of \([1,n]\) , since the \(E_{i}\) are bimodules for \(2 \leqslant i \leqslant n - 1\) . When n = 3, we have thus defined the Z-module \(E \otimes_{A} F \otimes_{B} G\) alluded to in no. 8, and recovered Proposition 8 (no. 8).

When each of the \(E_{i}\) is a multimodule (with, for \(2 \leqslant i \leqslant n - 1\) , \(A_{i-1}\) one of the rings operating on the left and \(A_{i}\) one of the rings operating on the right and analogous conditions for i = 1 and i = n), as in no. 4, a multimodule structure is defined on \(E_{1} \otimes_{A_1} E_{2} \otimes \cdots \otimes_{A_{n-1}} E_{n}\) with respect to all the rings except the \(A_{i}\) which operate on the \(E_{i}\) ( \(1 \leqslant i \leqslant n\) ).

In particular, let C be a commutative ring, \((\mathbf{E}_{i})_{1\leqslant i\leqslant n}\) a family of C-modules. By giving \(E_{1}\) and \(E_{n}\) two C-module structures identical with the given structure and \(E_{i}\) for \(2\leqslant i\leqslant n-1\) three C-module structures identical with the given structure, we define on the tensor product

\[
\mathbf {E} _ {1} \otimes_ {\mathbf {c}} \mathbf {E} _ {2} \otimes_ {\mathbf {c}} \mathbf {E} _ {3} \otimes \dots \otimes_ {\mathbf {c}} \mathbf {E} _ {n - 1} \otimes_ {\mathbf {c}} \mathbf {E} _ {n}\tag{34}
\]

\(n\) C-modules structures which are compatible with one another and which are in fact identical, since for all \(\gamma \in \mathbf{C}\) and \((x_{i})\in \prod_{i = 1}^{n}\mathbf{E}_{i}\) , by definition

\[
\begin{array}{r l} \left(\gamma x _ {1}\right) \otimes x _ {2} \otimes \dots \otimes x _ {n} \\ & = x _ {1} \otimes \left(\gamma x _ {2}\right) \otimes \dots \otimes x _ {n} = \dots = x _ {1} \otimes x _ {2} \otimes \dots \otimes \left(\gamma x _ {n}\right). \end{array}
\]

When we speak of the tensor product (34) as a C-module, we shall always mean with this structure, unless otherwise mentioned, and the tensor product (34) is also denoted by \(\bigotimes_{1\leqslant i\leqslant n}\mathbf{E}_i\) if no confusion can arise. For every C-module G, the Z-multilinear mappings of \(\prod_{i=1}^{n}\mathbf{E}_i\) into G which, for every index \(i\) , satisfy the relation

\[
f (x _ {1}, \dots , x _ {i - 1}, \gamma x _ {i}, x _ {i + 1}, \dots , x _ {n}) = \gamma f (x _ {1}, \dots , x _ {n})\tag{35}
\]

for \(\gamma \in \mathbf{C}\) and \((x_{i})\in \prod_{i}\mathbf{E}_{i}\) are then called C-multilinear and form a C-module denoted by \(\mathcal{L}_n(\mathbf{E}_1,\dots ,\mathbf{E}_n;\mathbf{G})\) ; the universal property of the tensor product (34) then allows us to define a canonical C-module isomorphism

\[
\mathcal {L} _ {n} \left(\mathrm{E} _ {1}, \dots , \mathrm{E} _ {n}; \mathrm{G}\right)\rightarrow \operatorname{Hom} _ {\mathrm{C}} \left(\mathrm{E} _ {1} \otimes_ {\mathrm{C}} \mathrm{E} _ {2} \otimes \dots \otimes_ {\mathrm{C}} \mathrm{E} _ {n}, \mathrm{G}\right)\tag{36}
\]

which associates with every C-multilinear mapping f the C-linear mapping g such that

\[
f (x _ {1}, \dots , x _ {n}) = g (x _ {1} \otimes x _ {2} \otimes \dots \otimes x _ {n}).
\]

A C-multilinear mapping of \(E_{1} \times \cdots \times E_{n}\) into C is also called an n-linear form.

Let \((\mathbf{F}_i)_{1\leqslant i\leqslant n}\) be another family of C-modules; for every system of \(n\) C-linear mappings \(u_i: \mathbf{E}_i \to \mathbf{F}_i\) , \(u_1 \otimes u_2 \otimes \cdots \otimes u_n\) (also denoted by \(\bigotimes_{1\leqslant i\leqslant n} u_i\) ) is a C-linear mapping of

\[
\mathbf {E} _ {1} \otimes_ {\mathbb {C}} \mathbf {E} _ {2} \otimes \dots \otimes_ {\mathbb {C}} \mathbf {E} _ {n} \quad \text { into } \mathbf {F} _ {1} \otimes_ {\mathbb {C}} \mathbf {F} _ {2} \otimes \dots \otimes_ {\mathbb {C}} \mathbf {F} _ {n}.
\]

Further, \((u_{1},\ldots,u_{n})\mapsto u_{1}\otimes u_{2}\otimes\cdots\otimes u_{n}\) is a C-multilinear mapping of \(\prod_{i}\operatorname{Hom}_{\mathbb{C}}(\mathbf{E}_{i},\mathbf{F}_{i})\) into

\[
\operatorname{Hom} _ {\mathbf {c}} (\mathbf {E} _ {1} \otimes_ {\mathbf {c}} \mathbf {E} _ {2} \otimes \dots \otimes_ {\mathbf {c}} \mathbf {E} _ {n}, \mathbf {F} _ {1} \otimes_ {\mathbf {c}} \mathbf {F} _ {2} \otimes \dots \otimes_ {\mathbf {c}} \mathbf {F} _ {n}).
\]

Hence there corresponds canonically to the latter mapping a C-linear mapping called canonical

\[
\begin{array}{r l} & \operatorname{Hom} _ {\mathbf {c}} (\mathrm{E} _ {1}, \mathrm{F} _ {1}) \otimes_ {\mathbf {c}} \operatorname{Hom} _ {\mathbf {c}} (\mathrm{E} _ {2}, \mathrm{F} _ {2}) \otimes \dots \otimes_ {\mathbf {c}} \operatorname{Hom} _ {\mathbf {c}} (\mathrm{E} _ {n}, \mathrm{F} _ {n}) \\ & \qquad \to \operatorname{Hom} _ {\mathbf {c}} (\mathrm{E} _ {1} \otimes_ {\mathbf {c}} \mathrm{E} _ {2} \otimes \dots \otimes_ {\mathbf {c}} \mathrm{E} _ {n}, \mathrm{F} _ {1} \otimes_ {\mathbf {c}} \mathrm{F} _ {2} \otimes \dots \otimes_ {\mathbf {c}} \mathrm{F} _ {n}) \end{array}\tag{37}
\]

generalizing that defined in no. 5 for \(n = 2\) .

The general associativity property seen earlier can be specialized here as follows. Given a partition \((\mathbf{J}_{k})_{1\leqslant k\leqslant m}\) of the interval \([1,n]\) of N, let, for each k, \(F_{k}\) be the tensor product \(E_{i_{1}}\otimes_{C}E_{i_{2}}\otimes\cdots\otimes_{C}E_{i_{r}}\) , where \((i_{1},\ldots,i_{r})\) is the strictly increasing sequence of elements of \(J_{k}\) . Then there is a canonical C-module isomorphism (called ``associativity isomorphism'')

\[
\mathrm{F} _ {1} \otimes_ {\mathbf {c}} \mathrm{F} _ {2} \otimes \dots \otimes_ {\mathbf {c}} \mathrm{F} _ {m} \rightarrow \mathrm{E} _ {1} \otimes_ {\mathbf {c}} \mathrm{E} _ {2} \otimes \dots \otimes_ {\mathbf {c}} \mathrm{E} _ {n}
\]

which, in the above notation, maps the tensor product

\[
y _ {1} \otimes y _ {2} \otimes \dots \otimes y _ {m}, \quad \text { where } y _ {k} = x _ {i _ {1}} \otimes x _ {i _ {2}} \otimes \dots \otimes x _ {i _ {r}},
\]

to the tensor product \(x_{1} \otimes x_{2} \otimes \cdots \otimes x_{n}\) (where \(x_{i} \in E_{i}\) for all i).

In particular, if \(\pi\) is a permutation of \([1, n]\) , writing \(\mathbf{J}_k = \{\pi(k)\}\) for \(1 \leqslant k \leqslant n\) , we obtain a canonical (``commutativity'') isomorphism

\[
\mathbf {E} _ {\pi (1)} \otimes_ {\mathbf {C}} \mathbf {E} _ {\pi (2)} \otimes \dots \otimes_ {\mathbf {C}} \mathbf {E} _ {\pi (n)} \rightarrow \mathbf {E} _ {1} \otimes_ {\mathbf {C}} \mathbf {E} _ {2} \otimes \dots \otimes_ {\mathbf {C}} \mathbf {E} _ {n}
\]

which maps \(x_{\pi(1)} \otimes x_{\pi(2)} \otimes \cdots \otimes x_{\pi(n)}\) to \(x_{1} \otimes x_{2} \otimes \cdots \otimes x_{n}\) . We shall often identify the various tensor products which correspond to one another under these canonical isomorphisms.

For \(1 \leqslant i \leqslant n\) , suppose that \(\mathbf{E}_i\) admits a basis \((b_{\lambda_1}^{(i)})_{\lambda_i \in \mathbf{L}_i}\) ; by induction on \(n\) , it follows from no. 7, Corollary 2 to Proposition 7 that the family \((b_{\lambda_1}^{(1)} \otimes b_{\lambda_2}^{(2)} \otimes \cdots \otimes b_{\lambda_n}^{(n)})\) , where \((\lambda_1, \ldots, \lambda_n)\) runs through \(\prod_{1 \leqslant i \leqslant n} \mathbf{L}_i\) , is a basis of \(\bigotimes_{1 \leqslant i \leqslant n} \mathbf{E}_i\) , sometimes called the tensor product of the bases \((b_{\lambda_1}^{(1)})\) in question.

Remarks. (1) The above remarks concerning the case of modules over a commutative ring generalize as in no. 5, Remark 2 when there is a tensor product \(\mathbf{E}_1 \otimes_{\mathbf{A}_1} \mathbf{E}_2 \otimes \cdots \otimes_{\mathbf{A}_{n-1}} \mathbf{E}_n\) where the rings \(\mathbf{A}_i\) are not necessarily commutative and where, for each \(i\) , there is a homomorphism \(\rho_i: \mathbf{C} \to \mathbf{A}_i\) of the same commutative ring \(\mathbf{C}\) such that: (1) \(\rho_i(\mathbf{C})\) is contained in the centre of \(\mathbf{A}_i\) ; (2) for \(2 \leqslant i \leqslant n - 1\) , the C-module structures on \(\mathbf{E}_i\) obtained using the homomorphisms \(\rho_{i-1}\) and \(\rho_i\) coincide. Then a C-module structure is obtained on \(\mathbf{E}_1 \otimes_{\mathbf{A}_1} \mathbf{E}_2 \otimes \cdots \otimes_{\mathbf{A}_{n-1}} \mathbf{E}_n\) and canonical mappings analogous to (13) (no. 5), which will be left to the reader to describe.

\begin{enumerate}
\def\labelenumi{(\arabic{enumi})}
\setcounter{enumi}{1}
\tightlist
\item
  Let A, B be two rings, E a right A-module, \(E'\) a left A-module, F a right B-module and \(F'\) a left B-module. The Z-bilinear mappings of \((\mathbf{E} \otimes_{\mathbf{A}} \mathbf{E}') \times (\mathbf{F} \otimes_{\mathbf{B}} \mathbf{F}')\) into a Z-module G are then in one-to-one correspondence with the Z-multilinear mappings f of \(E \times E' \times F \times F'\) into G satisfying the conditions
\end{enumerate}

\[
\left\{ \begin{array}{l} f (x \lambda , x ^ {\prime}, y, y ^ {\prime}) = f (x, \lambda x ^ {\prime}, y, y ^ {\prime}) \\ f (x, x ^ {\prime}, y \mu , y ^ {\prime}) = f (x, x ^ {\prime}, y, \mu y ^ {\prime}) \end{array} \right.\tag{38}
\]

for \(\lambda\in\mathbf{A}\) , \(\mu\in\mathbf{B}\) , \(x\in\mathbf{E}\) , \(x'\in\mathbf{E}'\) , \(y\in\mathbf{F}\) , \(y'\in\mathbf{F}'\) . The general constructions given in this no. reduce the proof of this to defining a canonical \(\mathbf{Z}\) -module isomorphism between \((\mathbf{E}\otimes_{\mathbf{A}}\mathbf{E}')\otimes_{\mathbf{Z}}(\mathbf{F}\otimes_{\mathbf{B}}\mathbf{F}')\) and \(\mathbf{E}\otimes_{\mathbf{A}}\mathbf{E}'\otimes_{\mathbf{Z}}\mathbf{F}\otimes_{\mathbf{B}}\mathbf{F}'\) , which follows from the associativity property of tensor products of the form (33).

